\bookchapter{第四十五章：質量保障——使AI所寫之代碼可靠}

\chapterintro{AI所寫之代碼有一個特點：它通常「看起來甚對」——變量命名規範、註釋齊全、邏輯清晰。但「看起來對」不等於「真之對」。AI代碼中隱藏之問題往往比手寫代碼更難發現，蓋其表面質量太高了，容易使人放鬆警惕。本章教汝如何為AI代碼建立有效之質量保障體系——用測試、審查、版本控制三道防線，將「看起來對」變成「真之對」。}

\booksection{一、測試驅動：使AI寫測試}

測試乃AI編程質量保障之第一道防線。好消息為：使AI寫測試比使它寫功能代碼更容易、更可靠——蓋測試之預期行為為明確之，不太容易產生「看起來對但實際不對」之代碼。

測試驅動之AI編程流程乃為：先使AI寫功能代碼，驗證基本功能通過後，立即使它寫測試用例。測試用例應覆蓋正常路徑、邊界情況與異常情況。然後再使AI根據測試反饋修正功能代碼。此流程之好處乃為：測試一旦寫好，即持續之驗證工具——後續任何修改皆可透過運行測試來確認無引入回歸問題。

\begin{lstlisting}
請爲 calculate_discount 函數編寫單元測試。
測試用例應覆蓋：
1. 正常路徑：
   - 100元商品打8折 -> 80
   - 200元商品打5折 -> 100
2. 邊界情況：
   - 0元商品打8折 -> 0
   - 100元商品打0折 -> 100（不打折）
   - 100元商品打10折 -> 100（全價）
3. 異常情況：
   - 負數價格 -> 應拋出錯誤
   - 折扣率 > 1 -> 應拋出錯誤
   - 非數字輸入 -> 應拋出錯誤
使用Vitest框架，予出完整測試文件。
\end{lstlisting}

注意此指令之特點：它明確列出了三類測試用例（正常、邊界、異常），不予AI自由發揮之空間去「選擇性測試」。AI有時候會只寫正常路徑之測試來「快速交差」，明確列出所有類別可避免此問題。

\booksection{二、代碼審查技巧}

即使有測試覆蓋，汝仍然須對AI產出之代碼進行人工審查。代碼審查不為逐行閱讀——那太耗時也太依賴汝之編程水平。有效之AI代碼審查當關注以下幾個關鍵維度：

\begin{center}
\begin{tabularx}{\linewidth}{|X|X|X|}
\hline
\textbf{審查維度} & \textbf{關注何} & \textbf{常見問題} \tabularnewline
\hline
安全性 & 輸入驗證、密鑰處理、SQL注入防護 & AI經常忘記驗證輸入，密鑰硬編碼於代碼中 \tabularnewline
錯誤處理 & 異常為否被捕獲，錯誤信息為否合理 & AI喜歡用空catch塊或只console.log \tabularnewline
性能 & 為否有明顯之性能問題，如N+1查詢、無限循環 & AI有時會寫出O(n²)之算法 \tabularnewline
可維護性 & 代碼結構為否清晰，命名為否有意義 & AI偶爾會生成過度抽象之代碼 \tabularnewline
依賴管理 & 引入之依賴為否必要，版本為否明確 & AI傾向於引入不必要之庫 \tabularnewline
\hline
\end{tabularx}
\end{center}

\begin{center}\small\color{gray}AI代碼審查之關鍵維度\end{center}

審查時不必自己逐一檢查——可使AI助汝做第一輪審查。指令如：「請審查此段代碼之安全性、錯誤處理與性能，指出潛於問題。」AI之自我審查能力於不斷提升，雖然不能完全替代人工審查，但能助汝過濾掉大部分明顯問題。

\booksection{三、版本控制：Git於AI編程中之角色}

Git於AI編程中之角色比於傳統編程中更加重要。蓋AI之每次修改皆可能引入不可預知之變化，汝須一個可靠之「後悔藥」——版本控制即那個後悔藥。

AI編程中Git之最佳實踐乃為「頻繁提交、小步快進」。每完成一個可驗證之功能點，就提交一次。如此當AI之某次修改出了問題時，汝可精準回退到上一個可用狀態，而不須推翻所有工作。

另一個實用技巧乃為「分支隔離」。對於結果不確定之AI修改（譬如大步重構），先創建一個新分支來做。若結果滿意，合併回主分支；若不滿意，直接刪除分支，不影響主分支。此比於主分支上反覆修改安全得多。

\begin{enumerate}
  \item 每次AI修改前，確保工作區為乾淨之（無未提交之更改）。如此若AI之修改不滿意，可用git checkout .快速回退。
  \item 每次AI修改後、驗證通過後，立即提交。提交信息描述此次修改做了何，譬如「feat: 添加用戶註冊功能」。
  \item 對於大步重構，先創建分支：git checkout -b refactor/auth-system。於分支上完成修改，測試通過後再合併。
  \item 養成查看diff之習慣。每次AI修改後，用git diff查看具體改了何。此不僅能助汝發現問題，還能助汝學習代碼。
\end{enumerate}

\booksection{四、常見陷阱與避坑指南}

AI編程中有幾個反覆出現之陷阱，幾乎所有初學者皆會踩。提前知道它等，能助汝少走甚多彎路。

\booksubsection{陷阱一：信任陷阱}

「AI所寫之代碼看起來甚專業，一定沒問題吧？」——此乃最危險之陷阱。AI生成之代碼於表面上甚整潔——命名規範、註釋齊全、結構清晰——此種「表面專業感」會使人放鬆警惕。但表面整潔不代表邏輯正確。解藥：永遠透過測試驗證，而非透過「看起來對不對」判斷。

\booksubsection{陷阱二：膨脹陷阱}

「既然AI能助吾寫代碼，那不如將功能做得全面一點？」——此會導致代碼膨脹。AI甚擅長「多予」：汝使它做一個登錄頁面，它會予汝一套完整之用戶管理系統。多餘之功能不僅增加複雜度，還增加了出bug之概率。解藥：嚴格遵循Simplicity First原則，明確告訴AI「只做此功能，勿額外添加」。

\booksubsection{陷阱三：理解陷阱}

「AI當理解吾之意思吧？」——不一定。AI之理解能力雖然甚強，但對含混之需求仍然會「猜」。當它猜錯了，產出之代碼即錯之。更麻煩者，它猜錯之代碼通常「看起來甚合理」，使汝誤以為它理解了。解藥：於指令中提供足夠之具體細節，減少AI須「猜」之空間。

\booksubsection{陷阱四：上下文陷阱}

「吾於之前之對話裡說過了，AI當還記得吧？」——甚可能不記得。隨著對話變長或開新會話，AI會丟失早期上下文。若某個關鍵約束無被寫入項目記憶檔案，它甚可能於後續修改中被忽略。解藥：關鍵約束寫入項目記憶檔案，勿依賴對話記憶。

\begin{tcbwarn}{四大陷阱速記}
信任陷阱——勿信「看起來對」；膨脹陷阱——AI愛「多予」；理解陷阱——含混需求它會「猜」；上下文陷阱——對話記憶會「忘」。對應之解藥分別為：測試驗證、砍掉多餘、予足細節、寫入記憶檔案。

\end{tcbwarn}

\begin{tcbpractice}{動手一試：為汝之項目建三道防線}
若汝已經開始做項目一（工具站）：(1) 使AI予核心函數寫3類測試用例；(2) 按五個維度審查一遍代碼（安全性／錯誤處理／性能／可維護性／依賴）；(3) 用git初始化項目，做第一次提交。三道防線齊了，汝之項目才真正屬於汝。

\end{tcbpractice}

\begin{tcbsummary}{本章小結}
\begin{itemize}
  \item AI代碼「看起來對」$\neq$「真之對」——表面質量太高反而容易放鬆警惕
  \item 使AI寫測試：明確列出正常／邊界／異常三類用例，防止「選擇性測試」
  \item 審查五維度：安全性、錯誤處理、性能、可維護性、依賴管理
  \item Git乃「後悔藥」：頻繁提交、小步快進、分支隔離、常看diff
  \item 四大陷阱：信任（測試驗證）、膨脹（砍掉多餘）、理解（予足細節）、上下文（寫入記憶檔案）
\end{itemize}

\end{tcbsummary}
